\documentclass[nojss,noheadings]{jss}

\usepackage{amsmath,amssymb,amsfonts}
\usepackage{bm}
\usepackage{booktabs}
\usepackage{tabularx}
\usepackage{enumitem}
\raggedbottom

\newcommand{\R}{\mathbb{R}}
\newcommand{\Rp}{\mathbb{R}_{+}}
\newcommand{\Ind}{\mathbb{I}}
\newcommand{\T}{\mathsf{T}}
\newcommand{\dd}{\mathrm{d}}
\renewcommand{\E}{\mathbb{E}}
\newcommand{\Var}{\mathbb{V}}
\newcommand{\Cov}{\mathbb{C}\mathrm{ov}}
\newcommand{\tr}{\operatorname{tr}}
\newcommand{\diag}{\operatorname{diag}}
\newcommand{\logit}{\operatorname{logit}}
\newcommand{\expit}{\operatorname{expit}}
\newcommand{\vect}[1]{\bm{#1}}
\newcommand{\mat}[1]{\mathbf{#1}}
\newcommand{\N}{\mathcal{N}}
\newcommand{\TNp}{\mathcal{N}^{+}}
\newcommand{\IG}{\operatorname{IG}}
\newcommand{\GIG}{\operatorname{GIG}}
\newcommand{\Exp}{\operatorname{Exp}}
\newcommand{\AL}{\operatorname{AL}}
\newcommand{\exAL}{\operatorname{exAL}}
\newcommand{\KL}{\operatorname{KL}}
\newcommand{\elbo}{\mathcal{L}}
\newcommand{\pigamma}{\pi_\gamma}
\renewcommand{\theequation}{S\arabic{equation}}
\renewcommand{\thetable}{S\arabic{table}}

\author{Antonio Aguirre \\ University of California Santa Cruz \And
        Raquel Barata \\ University of California Santa Cruz \AND
        Raquel Prado \\ University of California Santa Cruz \And
        Bruno Sans\'{o} \\ University of California Santa Cruz}
\title{Supplementary Theory and Algorithm Details for \pkg{exdqlm}}

\Plainauthor{Antonio Aguirre, Raquel Barata, Raquel Prado, Bruno Sanso}
\Plaintitle{Supplementary Theory and Algorithm Details for exdqlm}
\Shorttitle{Supplementary theory for \pkg{exdqlm}}

\Abstract{This supplement gives the mathematical details behind the inference
algorithms implemented in \pkg{exdqlm}. It records model hierarchies,
augmented posterior targets, Markov chain Monte Carlo (MCMC) updates,
variational Bayes (VB) factorizations, and evidence lower bound (ELBO)
components for dynamic quantile linear models (DQLMs) under asymmetric Laplace
(AL) and extended asymmetric Laplace (exAL) likelihoods, and for static AL/exAL
regression with ridge or regularized horseshoe normal-scale (RHS-NS) priors.
The notation is aligned with the main article and with the package interfaces:
vectors are written in bold lowercase, matrices in bold uppercase, and
dimensions are stated when objects are introduced.}

\Keywords{Bayesian quantile regression, dynamic linear models, asymmetric
Laplace likelihood, extended asymmetric Laplace likelihood, MCMC, variational
Bayes}
\Plainkeywords{Bayesian quantile regression, dynamic linear models, asymmetric
Laplace likelihood, extended asymmetric Laplace likelihood, MCMC, variational
Bayes}

\Address{
  Antonio Aguirre and Raquel Barata\\
  Department of Statistics\\
  University of California Santa Cruz\\
  Santa Cruz, CA 95064\\
  E-mail: \email{jaguir26@ucsc.edu}, \email{raquel.a.barata@gmail.com}\\
}

\begin{document}

\section{Scope and organization}

This supplement is a mathematical companion to the main article. It is intended
for readers who want to inspect the posterior targets and update equations
behind the package routines, rather than for users who only need the applied
workflow shown in the examples. The dynamic sections cover AL-based DQLMs and
exAL-based exDQLMs. The static sections cover AL/exAL regression with Gaussian
ridge priors and with the RHS-NS sparse prior used by the package.

Each model section follows the same organization: model hierarchy, augmented
posterior target, MCMC updates, VB factorization and coordinate updates, ELBO
components, and the corresponding package function. The final section maps the
main user-facing functions to these algorithmic blocks.

\section{Notation and distributional building blocks}

Let $p_0\in(0,1)$ be the target quantile. The response is scalar. In dynamic
models, observations are indexed by $t=1,\ldots,T$; in static regression,
observations are indexed by $i=1,\ldots,n$. Lowercase bold symbols denote
vectors, for example $\vect{\theta}_t\in\R^q$ and $\vect{\beta}\in\R^p$.
Uppercase bold symbols denote matrices, for example $\mat{G}_t\in\R^{q\times q}$
and $\mat{X}\in\R^{n\times p}$. For any positive definite matrix $\mat{A}$,
$|\mat{A}|$ denotes its determinant.

\subsection{Notation contract}

The supplement uses the following conventions throughout. The symbols
$p_\gamma$, $A_\gamma$, $B_\gamma$, and $C_\gamma$ are deterministic functions
of $(p_0,\gamma)$, whereas $\pigamma(\gamma)$ denotes the prior density for
$\gamma$ on $(L,U)$. Generic densities are denoted by $p(\cdot)$ only when no
confusion with the exAL quantile map can arise. In variational sections,
expectations without an explicit subscript are taken under the current
mean-field distribution $q$. The symbol $\xi$ is reserved for the RHS-NS
half-Cauchy auxiliary variable; the transformed coordinate for $\gamma$ in
Laplace-Delta calculations is denoted by $z_\gamma$.

\begin{table}[!htb]
\centering
\begin{tabularx}{\textwidth}{@{}lX@{}}
\toprule
Symbol & Meaning \\
\midrule
$p_0$ & Target quantile level. \\
$p_\gamma$ & exAL internal quantile map determined by $(p_0,\gamma)$. \\
$A_0,B_0$ & AL constants $A(p_0)$ and $B(p_0)$. \\
$A_\gamma,B_\gamma,C_\gamma$ & exAL constants evaluated at $p_\gamma$. \\
$\pigamma(\gamma)$ & Prior density for $\gamma$ on $(L,U)$. \\
$\vect{\theta}_t,\vect{f}_t,\mat{G}_t,\mat{W}_t$ & Dynamic state, observation vector, evolution matrix, and evolution covariance. \\
$\vect{\beta},\mat{X}$ & Static regression coefficients and design matrix. \\
$\mathcal{A}$ & Active coefficient set receiving RHS-NS shrinkage. \\
\bottomrule
\end{tabularx}
\caption{Core notation used throughout the supplement.}
\label{tab:supp_notation}
\end{table}

The distributional parameterizations are also fixed throughout:
$\N(m,V)$ uses variance $V$; $\IG(a,b)$ uses rate $b$ with density proportional
to $x^{-(a+1)}\exp(-b/x)$; $\Exp(\text{rate}=1/\sigma)$ has density
$\sigma^{-1}\exp(-x/\sigma)$; and $\TNp(m,V)$ denotes $\N(m,V)$ truncated to
$(0,\infty)$. The generalized inverse Gaussian parameterization is given in
\eqref{eq:gig_parameterization}.

The asymmetric Laplace (AL) mixture representation used by the package is
\begin{align}
v &\mid \sigma \sim \Exp(\text{rate}=1/\sigma), \qquad v>0, \\
y \mid \eta,\sigma,v
&\sim \N\{ \eta + A(p_0)v,\ \sigma B(p_0)v \},
\end{align}
where
\begin{equation}
A(p)=\frac{1-2p}{p(1-p)}, \qquad B(p)=\frac{2}{p(1-p)}.
\label{eq:al_AB}
\end{equation}
The AL model is the $\gamma=0$ special case of the exAL formulation below.

The exAL distribution is represented through an additional latent variable
$s\sim\TNp(0,1)$, a standard normal truncated to $(0,\infty)$. For
$\gamma\in(L,U)$, define
\begin{align}
g(\gamma) &= 2\Phi(-|\gamma|)\exp(\gamma^2/2), \\
p_\gamma &= \Ind(\gamma<0) + \frac{p_0-\Ind(\gamma<0)}{g(\gamma)}, \\
A_\gamma &= A(p_\gamma), \qquad
B_\gamma = B(p_\gamma), \qquad
C_\gamma = \{\Ind(\gamma>0)-p_\gamma\}^{-1}.
\label{eq:exal_ABC}
\end{align}
The endpoints $L<0<U$ are the roots $g(L)=1-p_0$ and $g(U)=p_0$. Conditional on
$(v,s,\sigma,\gamma)$,
\begin{equation}
y\mid \eta,\sigma,\gamma,v,s
\sim
\N\left(\eta + C_\gamma\sigma|\gamma|s + A_\gamma v,\ 
\sigma B_\gamma v \right).
\label{eq:exal_aug}
\end{equation}
The package uses an inverse-gamma prior $\sigma\sim\IG(a_\sigma,b_\sigma)$
and a proper truncated Student-$t$ prior $\pigamma(\gamma)$ for $\gamma$ on
$(L,U)$.

The generalized inverse Gaussian distribution is parameterized as
\begin{equation}
x\sim\GIG(\lambda,\chi,\psi)
\quad\Longleftrightarrow\quad
p(x)\propto x^{\lambda-1}
\exp\left\{-\frac{1}{2}\left(\frac{\chi}{x}+\psi x\right)\right\},
\qquad x>0.
\label{eq:gig_parameterization}
\end{equation}
For $r\in\R$,
\begin{equation}
\E(x^r)=
\left(\frac{\chi}{\psi}\right)^{r/2}
\frac{K_{\lambda+r}(\sqrt{\chi\psi})}{K_{\lambda}(\sqrt{\chi\psi})},
\label{eq:gig_moments}
\end{equation}
where $K_\nu(\cdot)$ is the modified Bessel function of the second kind.

\section{Dynamic DQLM under the AL likelihood}
\label{sec:dyn_dqlm}

\subsection{Model hierarchy and full joint}

Let $\vect{\theta}_t\in\R^q$ be the latent state, $\vect{f}_t\in\R^q$ the
observation vector, $\mat{G}_t\in\R^{q\times q}$ the evolution matrix, and
$\mat{W}_t\in\R^{q\times q}$ a positive semidefinite evolution covariance. The
dynamic DQLM is
\begin{align}
\vect{\theta}_0 &\sim \N(\vect{m}_0,\mat{C}_0), \qquad
  \vect{m}_0\in\R^q,\ \mat{C}_0\in\R^{q\times q}, \\
\vect{\theta}_t\mid\vect{\theta}_{t-1}
&\sim \N(\mat{G}_t\vect{\theta}_{t-1},\mat{W}_t), \\
v_t\mid\sigma &\sim \Exp(\text{rate}=1/\sigma), \\
y_t\mid\vect{\theta}_t,\sigma,v_t
&\sim \N(\vect{f}_t^\T\vect{\theta}_t + A_0 v_t,\ \sigma B_0 v_t),
\label{eq:dyn_dqlm_hierarchy}
\end{align}
where $A_0=A(p_0)$ and $B_0=B(p_0)$. With
$\vect{v}=(v_1,\ldots,v_T)^\T$ and
$\vect{\theta}_{0:T}=(\vect{\theta}_0,\ldots,\vect{\theta}_T)$, the posterior
target is
\begin{align}
p(\vect{\theta}_{0:T},\sigma,\vect{v}\mid\vect{y})
&\propto
p(\sigma)\,p(\vect{\theta}_0)
\prod_{t=1}^T p(\vect{\theta}_t\mid\vect{\theta}_{t-1})
\prod_{t=1}^T p(v_t\mid\sigma)
\prod_{t=1}^T p(y_t\mid\vect{\theta}_t,\sigma,v_t).
\label{eq:dyn_dqlm_joint}
\end{align}
Equivalently, conditional on $(\sigma,\vect{v})$, the pseudo-observation
$\tilde y_t=y_t-A_0v_t$ satisfies
\begin{equation}
\tilde y_t=\vect{f}_t^\T\vect{\theta}_t+\epsilon_t,\qquad
\epsilon_t\sim\N(0,\sigma B_0v_t).
\label{eq:dyn_dqlm_pseudo}
\end{equation}
Thus the conditional state posterior is a Gaussian state-space posterior.

\subsection{MCMC updates}

One Gibbs sweep for the dynamic DQLM uses the following full conditionals.

\paragraph{State path.}
Given $(\sigma,\vect{v})$, sample the state path
$\vect{\theta}_{0:T}$ exactly by applying forward-filtering
backward-sampling (FFBS) to \eqref{eq:dyn_dqlm_pseudo}.

\paragraph{Scale.}
Let $r_t=y_t-\vect{f}_t^\T\vect{\theta}_t$. Then
\begin{equation}
\sigma\mid \vect{\theta}_{0:T},\vect{v},\vect{y}
\sim \IG\left(
a_\sigma+\frac{3T}{2},\
b_\sigma+\sum_{t=1}^T v_t+
\sum_{t=1}^T\frac{(r_t-A_0v_t)^2}{2B_0v_t}
\right).
\label{eq:dyn_dqlm_sigma_cond}
\end{equation}

\paragraph{Latent AL variables.}
For $t=1,\ldots,T$,
\begin{equation}
v_t\mid\vect{\theta}_t,\sigma,y_t
\sim
\GIG\left(
\frac{1}{2},\
\frac{r_t^2}{\sigma B_0},\
\frac{A_0^2}{\sigma B_0}+\frac{2}{\sigma}
\right).
\label{eq:dyn_dqlm_v_cond}
\end{equation}

\subsection{Mean-field VB and ELBO}

The dynamic DQLM VB approximation uses
\begin{equation}
q(\vect{\theta}_{0:T},\sigma,\vect{v})
=q(\vect{\theta}_{0:T})q(\sigma)\prod_{t=1}^Tq(v_t),
\label{eq:dyn_dqlm_mf}
\end{equation}
with $q(\sigma)=\IG(a_q,b_q)$ and
$q(v_t)=\GIG(1/2,\chi_t,\psi_t)$. Define
\begin{equation}
\kappa=\E_q(\sigma^{-1})=\frac{a_q}{b_q},\qquad
\ell_t=\E_q(v_t^{-1}),\qquad
\nu_t=\E_q(v_t).
\end{equation}
The state factor is the Gaussian posterior of the auxiliary model
\begin{equation}
y_t^\star=\vect{f}_t^\T\vect{\theta}_t+\epsilon_t^\star,\qquad
\epsilon_t^\star\sim\N(0,R_t^\star),
\qquad
y_t^\star=y_t-\frac{A_0}{\ell_t},\quad
R_t^\star=\frac{B_0}{\kappa\ell_t}.
\label{eq:dyn_dqlm_vb_pseudo}
\end{equation}
After Kalman smoothing, write
$q(\vect{\theta}_t)=\N(\vect{m}_{t|T},\mat{C}_{t|T})$ and define
\begin{align}
\bar r_t &= y_t-\vect{f}_t^\T\vect{m}_{t|T},\\
\overline{r_t^2}
&= \bar r_t^2+\vect{f}_t^\T\mat{C}_{t|T}\vect{f}_t.
\end{align}
The remaining coordinate-ascent variational inference (CAVI) updates are
\begin{align}
q(v_t)&=\GIG\left(
\frac{1}{2},\
\frac{\kappa}{B_0}\overline{r_t^2},\
\kappa\left(2+\frac{A_0^2}{B_0}\right)\right), \\
q(\sigma)&=\IG(a_q,b_q), \qquad
a_q=a_\sigma+\frac{3T}{2},\\
b_q&=b_\sigma+\sum_{t=1}^T\nu_t+
\frac{1}{2B_0}\sum_{t=1}^T
\left\{\ell_t\overline{r_t^2}-2A_0\bar r_t+A_0^2\nu_t\right\}.
\label{eq:dyn_dqlm_vb_sigma}
\end{align}

The ELBO is
\begin{align}
\elbo(q)
&=
\E_q\{\log p(\vect{y},\vect{\theta}_{0:T},\sigma,\vect{v})\}
-\E_q\{\log q(\vect{\theta}_{0:T},\sigma,\vect{v})\} \\
&=\elbo_\theta+\E_q\{\log p(\sigma)\}
+\sum_{t=1}^T\E_q\{\log p(v_t\mid\sigma)\}
+\sum_{t=1}^T\E_q\{\log p(y_t\mid\vect{\theta}_t,\sigma,v_t)\}\\
&\quad
+H\{q(\sigma)\}+\sum_{t=1}^TH\{q(v_t)\}.
\label{eq:dyn_dqlm_elbo}
\end{align}
The state-path term $\elbo_\theta$ is evaluated from the auxiliary Gaussian
state-space model in \eqref{eq:dyn_dqlm_vb_pseudo} through the Kalman marginal
likelihood identity used in the package ELBO diagnostics.

\section{Dynamic exDQLM under the exAL likelihood}
\label{sec:dyn_exdqlm}

\subsection{Model hierarchy and full joint}

The dynamic exDQLM replaces the AL observation in
\eqref{eq:dyn_dqlm_hierarchy} with the exAL augmentation
\begin{align}
v_t\mid\sigma &\sim \Exp(\text{rate}=1/\sigma), \qquad
s_t\sim\TNp(0,1), \\
y_t\mid\vect{\theta}_t,\sigma,\gamma,v_t,s_t
&\sim
\N\left(
\vect{f}_t^\T\vect{\theta}_t+C_\gamma\sigma|\gamma|s_t+A_\gamma v_t,\
\sigma B_\gamma v_t
\right),
\label{eq:dyn_exdqlm_obs}
\end{align}
with the same state evolution as in Section~\ref{sec:dyn_dqlm}. The complete
posterior target is
\begin{align}
p(\vect{\theta}_{0:T},\vect{v},\vect{s},\sigma,\gamma\mid\vect{y})
&\propto
p(\vect{\theta}_0)\prod_{t=1}^Tp(\vect{\theta}_t\mid\vect{\theta}_{t-1})
\prod_{t=1}^Tp(y_t\mid\vect{\theta}_t,v_t,s_t,\sigma,\gamma) \nonumber\\
&\quad\times
\prod_{t=1}^Tp(v_t\mid\sigma)\prod_{t=1}^Tp(s_t)\,p(\sigma)\pigamma(\gamma).
\label{eq:dyn_exdqlm_joint}
\end{align}
Conditional on $(\vect{v},\vect{s},\sigma,\gamma)$,
\begin{equation}
\tilde y_t=y_t-A_\gamma v_t-C_\gamma\sigma|\gamma|s_t,\qquad
R_t=\sigma B_\gamma v_t,
\label{eq:dyn_exdqlm_pseudo}
\end{equation}
which again gives a linear Gaussian state-space model for
$\vect{\theta}_{0:T}$.

\subsection{MCMC updates}

The default package path samples $\sigma$ from its exact conditional
given $\gamma$ and then updates $\gamma$ by bounded univariate slice sampling.
Joint random-walk alternatives are also implemented, but the following target
is the default exact-kernel path.

\paragraph{Latent $v_t$.}
Let
$r_t=y_t-\vect{f}_t^\T\vect{\theta}_t-C_\gamma\sigma|\gamma|s_t$. Then
\begin{equation}
v_t\mid\cdot
\sim
\GIG\left(
\frac{1}{2},\
\frac{r_t^2}{\sigma B_\gamma},\
\frac{A_\gamma^2}{\sigma B_\gamma}+\frac{2}{\sigma}
\right).
\label{eq:dyn_exdqlm_v_cond}
\end{equation}

\paragraph{Latent $s_t$.}
Let $y_t^\circ=y_t-\vect{f}_t^\T\vect{\theta}_t-A_\gamma v_t$ and
$d_\gamma=C_\gamma\sigma|\gamma|$. Then
\begin{equation}
s_t\mid\cdot\sim\TNp(m_{s,t},V_{s,t}),
\qquad
V_{s,t}=\left(1+\frac{d_\gamma^2}{\sigma B_\gamma v_t}\right)^{-1},
\qquad
m_{s,t}=V_{s,t}\frac{d_\gamma y_t^\circ}{\sigma B_\gamma v_t}.
\label{eq:dyn_exdqlm_s_cond}
\end{equation}

\paragraph{State path.}
Given $(\vect{v},\vect{s},\sigma,\gamma)$, sample
$\vect{\theta}_{0:T}$ by FFBS using \eqref{eq:dyn_exdqlm_pseudo}.

\paragraph{Scale.}
For fixed $\gamma$, let $y_t^\circ=y_t-\vect{f}_t^\T\vect{\theta}_t-A_\gamma v_t$
and $u_t=C_\gamma|\gamma|s_t$. Then
\begin{equation}
\sigma\mid\cdot
\sim
\GIG(\lambda_\sigma,\chi_\sigma,\psi_\sigma),
\label{eq:dyn_exdqlm_sigma_cond}
\end{equation}
where
\begin{align}
\lambda_\sigma &= -\left(a_\sigma+\frac{3T}{2}\right), \\
\chi_\sigma &= 2b_\sigma+2\sum_{t=1}^T v_t+
\sum_{t=1}^T\frac{(y_t^\circ)^2}{B_\gamma v_t},\\
\psi_\sigma &= \sum_{t=1}^T\frac{u_t^2}{B_\gamma v_t}.
\end{align}

\paragraph{Skewness.}
The conditional kernel for $\gamma$ is
\begin{align}
\pi(\gamma\mid\cdot)
&\propto \pigamma(\gamma)\,
\prod_{t=1}^T(\sigma B_\gamma v_t)^{-1/2}
\exp\left[
-\sum_{t=1}^T
\frac{
\{y_t-\vect{f}_t^\T\vect{\theta}_t-A_\gamma v_t-C_\gamma\sigma|\gamma|s_t\}^2
}{2\sigma B_\gamma v_t}
\right],
\label{eq:dyn_exdqlm_gamma_kernel}
\end{align}
for $\gamma\in(L,U)$. The default sampler applies the bounded slice update in
Section~\ref{sec:slice} directly to this kernel.

\subsection{Laplace-Delta VB approximation}

The mean-field family is
\begin{equation}
q(\vect{\theta}_{0:T},\vect{v},\vect{s},\sigma,\gamma)
=q(\vect{\theta}_{0:T})\prod_{t=1}^Tq(v_t)\prod_{t=1}^Tq(s_t)\,q(\sigma,\gamma).
\label{eq:dyn_exdqlm_mf}
\end{equation}
The conjugate factors are updated conditionally on moments from
$q(\sigma,\gamma)$. Define
\begin{align}
\kappa_1 &= \E\{1/(\sigma B_\gamma)\},&
\kappa_2 &= \E\{A_\gamma/(\sigma B_\gamma)\},&
\kappa_3 &= \E\{A_\gamma^2/(\sigma B_\gamma)\},\\
\kappa_4 &= \E\{C_\gamma|\gamma|/B_\gamma\},&
\kappa_5 &= \E\{\sigma C_\gamma^2\gamma^2/B_\gamma\},&
\kappa_6 &= \E\{A_\gamma C_\gamma|\gamma|/B_\gamma\}.
\label{eq:dyn_exdqlm_kappa}
\end{align}
Let $m_{1/v,t}=\E(v_t^{-1})$, $m_{s,t}=\E(s_t)$, and
$m_{s^2,t}=\E(s_t^2)$. The state factor is Gaussian and is computed by Kalman
smoothing with observation precision
\begin{equation}
\phi_t=\kappa_1m_{1/v,t}
\end{equation}
and pseudo-response
\begin{equation}
y_t^{\mathrm{VB}}
=
\frac{
y_t\phi_t-\kappa_2-\kappa_4m_{s,t}m_{1/v,t}
}{\phi_t}.
\label{eq:dyn_exdqlm_vb_pseudo}
\end{equation}
If $\eta_t=\vect{f}_t^\T\vect{\theta}_t$, write
$m_{\eta,t}=\E(\eta_t)$, $S_{\eta,t}=\Var(\eta_t)$,
$m_{\Delta,t}=y_t-m_{\eta,t}$, and
$S_{\Delta,t}=m_{\Delta,t}^2+S_{\eta,t}$. Then
\begin{align}
q(v_t)&=\GIG\left(
\frac{1}{2},\
\kappa_1S_{\Delta,t}-2\kappa_4m_{\Delta,t}m_{s,t}+\kappa_5m_{s^2,t},\
\kappa_3+2\E(\sigma^{-1})
\right),\\
q(s_t)&=\TNp(m_{s,t}^{\mathrm{VB}},V_{s,t}^{\mathrm{VB}}),\\
V_{s,t}^{\mathrm{VB}}
&=\left(1+\kappa_5m_{1/v,t}\right)^{-1},\\
m_{s,t}^{\mathrm{VB}}
&=V_{s,t}^{\mathrm{VB}}
\left(\kappa_4m_{\Delta,t}m_{1/v,t}-\kappa_6\right).
\label{eq:dyn_exdqlm_vb_vs}
\end{align}
The nonconjugate factor $q(\sigma,\gamma)$ is approximated by a
Laplace-Delta step on transformed coordinates
\begin{equation}
u=\log\sigma,\qquad
z_\gamma=\logit\left(\frac{\gamma-L}{U-L}\right).
\label{eq:dyn_exdqlm_ld_transform}
\end{equation}
The transformed objective is the expected complete log-joint plus the
Jacobian
\begin{equation}
u+\log\{\dd\gamma/\dd z_\gamma\}.
\end{equation}
The package stores the transformed mode, the approximate covariance matrix,
Delta-method moments, and ELBO components.

\paragraph{ELBO.}
The LDVB ELBO has the block decomposition
\begin{equation}
\elbo(q)=
\elbo_{\mathrm{like}}+\elbo_\theta+\elbo_v+\elbo_s+\elbo_{\sigma\gamma}
+H\{q(\vect{\theta}_{0:T})\}
+\sum_{t=1}^TH\{q(v_t)\}
+\sum_{t=1}^TH\{q(s_t)\}
+H\{q(\sigma,\gamma)\}.
\label{eq:dyn_exdqlm_elbo}
\end{equation}
The state and latent-variable blocks are exact for the current
$q(\sigma,\gamma)$ moments; the $(\sigma,\gamma)$ entropy and smooth
expectations are evaluated under the Laplace-Delta approximation.

\section{Static AL and exAL regression with ridge priors}
\label{sec:static_ridge}

\subsection{Model hierarchy and full joint}

Let $\mat{X}\in\R^{n\times p}$ have row vectors $\vect{x}_i^\T$, and let
$\vect{\beta}\in\R^p$. The ridge/Gaussian coefficient prior is
\begin{equation}
\vect{\beta}\sim\N(\vect{b}_0,\mat{B}_0),\qquad
\vect{b}_0\in\R^p,\quad \mat{B}_0\in\R^{p\times p}.
\label{eq:static_beta_ridge}
\end{equation}
Under the exAL likelihood,
\begin{align}
v_i\mid\sigma &\sim \Exp(\text{rate}=1/\sigma),\qquad
s_i\sim\TNp(0,1),\\
y_i\mid\vect{\beta},\sigma,\gamma,v_i,s_i
&\sim
\N\left(
\vect{x}_i^\T\vect{\beta}+C_\gamma\sigma|\gamma|s_i+A_\gamma v_i,\
\sigma B_\gamma v_i
\right).
\label{eq:static_exal_hierarchy}
\end{align}
The static AL model is obtained by setting $\gamma=0$, removing $s_i$, and
using $A_0,B_0$. The exAL full joint is
\begin{align}
p(\vect{\beta},\sigma,\gamma,\vect{v},\vect{s},\vect{y})
&=
p(\vect{\beta})p(\sigma)\pigamma(\gamma)
\prod_{i=1}^n
p(y_i\mid\vect{\beta},\sigma,\gamma,v_i,s_i)
p(v_i\mid\sigma)p(s_i).
\label{eq:static_exal_joint}
\end{align}

\subsection{MCMC updates}

For fixed $(\sigma,\gamma,\vect{v},\vect{s})$, define
\begin{equation}
y_i^\star=y_i-C_\gamma\sigma|\gamma|s_i-A_\gamma v_i,\qquad
w_i=(\sigma B_\gamma v_i)^{-1}.
\end{equation}
Let $\mat{W}=\diag(w_1,\ldots,w_n)$ and
$\vect{y}^\star=(y_1^\star,\ldots,y_n^\star)^\T$. Then
\begin{align}
\mat{\Sigma}_{\beta\mid\cdot}
&=\left(\mat{B}_0^{-1}+\mat{X}^\T\mat{W}\mat{X}\right)^{-1},\\
\vect{m}_{\beta\mid\cdot}
&=\mat{\Sigma}_{\beta\mid\cdot}
\left(\mat{B}_0^{-1}\vect{b}_0+\mat{X}^\T\mat{W}\vect{y}^\star\right),\\
\vect{\beta}\mid\cdot
&\sim\N(\vect{m}_{\beta\mid\cdot},\mat{\Sigma}_{\beta\mid\cdot}).
\label{eq:static_beta_cond}
\end{align}
Let $D_\gamma=C_\gamma|\gamma|$. The remaining exAL full conditionals are
\begin{align}
s_i\mid\cdot
&\sim \TNp(m_{s,i},V_{s,i}),\\
V_{s,i}
&=\left(1+\frac{\sigma D_\gamma^2}{B_\gamma v_i}\right)^{-1},\\
m_{s,i}
&=V_{s,i}\frac{D_\gamma(y_i-\vect{x}_i^\T\vect{\beta}-A_\gamma v_i)}
{B_\gamma v_i},
\label{eq:static_exal_s_cond}
\end{align}
and, with $r_i=y_i-\vect{x}_i^\T\vect{\beta}-\sigma D_\gamma s_i$,
\begin{equation}
v_i\mid\cdot
\sim
\GIG\left(
\frac{1}{2},\
\frac{r_i^2}{\sigma B_\gamma},\
\frac{A_\gamma^2}{\sigma B_\gamma}+\frac{2}{\sigma}
\right).
\label{eq:static_exal_v_cond}
\end{equation}
For fixed $\gamma$, define $t_i=y_i-\vect{x}_i^\T\vect{\beta}-A_\gamma v_i$.
Then
\begin{equation}
\sigma\mid\cdot
\sim
\GIG\left(
-a_\sigma-\frac{3n}{2},\
2b_\sigma+2\sum_{i=1}^n v_i+\sum_{i=1}^n\frac{t_i^2}{B_\gamma v_i},\
\sum_{i=1}^n\frac{D_\gamma^2s_i^2}{B_\gamma v_i}
\right).
\label{eq:static_exal_sigma_cond}
\end{equation}
The nonstandard $\gamma$ kernel is
\begin{align}
\pi(\gamma\mid\cdot)
&\propto
\pigamma(\gamma)\,[B_\gamma]^{-n/2}
\exp\left\{
\sum_{i=1}^n\frac{D_\gamma s_it_i}{B_\gamma v_i}
-\frac{1}{2\sigma}\sum_{i=1}^n\frac{t_i^2}{B_\gamma v_i}
-\frac{\sigma}{2}\sum_{i=1}^n\frac{D_\gamma^2s_i^2}{B_\gamma v_i}
\right\},
\label{eq:static_exal_gamma_cond}
\end{align}
on $(L,U)$. The package default updates this one-dimensional target by the
bounded slice sampler in Section~\ref{sec:slice}.

For the AL special case, the Gibbs updates reduce to
\begin{align}
\vect{\beta}\mid\cdot
&\sim\N(\vect{m}_{\beta\mid\cdot},\mat{\Sigma}_{\beta\mid\cdot}),\\
\sigma\mid\cdot
&\sim
\IG\left(
a_\sigma+\frac{3n}{2},\
b_\sigma+\sum_{i=1}^n v_i+
\sum_{i=1}^n\frac{(y_i-\vect{x}_i^\T\vect{\beta}-A_0v_i)^2}{2B_0v_i}
\right),\\
v_i\mid\cdot
&\sim
\GIG\left(
\frac{1}{2},\
\frac{(y_i-\vect{x}_i^\T\vect{\beta})^2}{\sigma B_0},\
\frac{A_0^2}{\sigma B_0}+\frac{2}{\sigma}
\right).
\label{eq:static_al_mcmc}
\end{align}

\subsection{Laplace-Delta VB updates and ELBO}

The static exAL variational family is
\begin{equation}
q(\vect{\beta},\vect{v},\vect{s},\sigma,\gamma)
=q(\vect{\beta})\prod_{i=1}^nq(v_i)\prod_{i=1}^nq(s_i)q(\sigma,\gamma).
\label{eq:static_exal_mf}
\end{equation}
The coefficient factor is Gaussian. With the same moment definitions as in
\eqref{eq:dyn_exdqlm_kappa},
\begin{align}
\omega_i&=\kappa_1\E(v_i^{-1}),\\
\eta_i&=
y_i\kappa_1\E(v_i^{-1})
-\kappa_4\E(v_i^{-1})\E(s_i)
-\kappa_2,\\
\mat{\Sigma}_\beta
&=\left(\mat{B}_0^{-1}+\mat{X}^\T\diag(\omega_1,\ldots,\omega_n)\mat{X}\right)^{-1},\\
\vect{m}_\beta
&=\mat{\Sigma}_\beta
\left(\mat{B}_0^{-1}\vect{b}_0+\mat{X}^\T\vect{\eta}\right),
\label{eq:static_exal_vb_beta}
\end{align}
where $\vect{\eta}=(\eta_1,\ldots,\eta_n)^\T$. The $q(v_i)$ and $q(s_i)$
updates are also closed form conditional on moments of $q(\sigma,\gamma)$.
Let
\begin{equation}
\bar r_i=y_i-\vect{x}_i^\T\vect{m}_\beta,\qquad
\overline{r_i^2}=\bar r_i^2+\vect{x}_i^\T\mat{\Sigma}_\beta\vect{x}_i.
\end{equation}
Then
\begin{align}
q(v_i)
&=
\GIG\left(
\frac{1}{2},\
\kappa_1\overline{r_i^2}
-2\kappa_4\bar r_i\E(s_i)
+\kappa_5\E(s_i^2),\
\kappa_3+2\E(\sigma^{-1})
\right),\\
q(s_i)&=\TNp(m_{s,i}^{\mathrm{VB}},V_{s,i}^{\mathrm{VB}}),\\
V_{s,i}^{\mathrm{VB}}
&=\left(1+\kappa_5\E(v_i^{-1})\right)^{-1},\\
m_{s,i}^{\mathrm{VB}}
&=V_{s,i}^{\mathrm{VB}}
\left\{\kappa_4\bar r_i\E(v_i^{-1})-\kappa_6\right\}.
\label{eq:static_exal_vb_vs}
\end{align}
The nonconjugate factor $q(\sigma,\gamma)$ is proportional to
\begin{align}
q(\sigma,\gamma)
&\propto
\pigamma(\gamma)[B_\gamma]^{-n/2}
\sigma^{-(a_\sigma+1+3n/2)}
\exp\left\{
-\frac{1}{2}\left(\frac{\bar c(\gamma)}{\sigma}
+\bar d(\gamma)\sigma\right)+\bar e(\gamma)
\right\},
\label{eq:static_exal_q_sigmagam}
\end{align}
where
\begin{align}
\bar c(\gamma)
&=2b_\sigma+2\sum_{i=1}^n\E(v_i)
+\sum_{i=1}^n
\frac{\E\{(y_i-\vect{x}_i^\T\vect{\beta}-A_\gamma v_i)^2/v_i\}}
{B_\gamma},\\
\bar d(\gamma)
&=\sum_{i=1}^n\frac{C_\gamma^2\gamma^2}{B_\gamma}\E\{s_i^2/v_i\},\\
\bar e(\gamma)
&=\sum_{i=1}^n
\frac{C_\gamma|\gamma|}{B_\gamma}
\E\{s_i(y_i-\vect{x}_i^\T\vect{\beta}-A_\gamma v_i)/v_i\}.
\end{align}
This factor is approximated by the same Laplace-Delta method as in the
dynamic exDQLM. The AL path sets $\gamma=0$ and uses
$q(\vect{\beta})q(\sigma)\prod_iq(v_i)$ with the same closed-form updates as
Section~\ref{sec:dyn_dqlm}, replacing state smoothing moments by
$\E\{(y_i-\vect{x}_i^\T\vect{\beta})^2\}$.

The static exAL ELBO is
\begin{equation}
\elbo(q)=
\E_q\{\log p(\vect{\beta})+\log p(\sigma)+\log \pigamma(\gamma)\}
+\sum_{i=1}^n\E_q\{\log p(y_i\mid\cdot)+\log p(v_i\mid\sigma)+\log p(s_i)\}
+H(q),
\label{eq:static_exal_elbo}
\end{equation}
where $H(q)$ is the sum of the Gaussian, GIG, truncated-normal, and
Laplace-Delta entropy terms.

\section{Static regression with the RHS-NS prior}
\label{sec:static_rhs}

This section focuses on the \code{rhs_ns} coefficient prior, the
Gibbs-friendly regularized horseshoe variant used for sparse static
regression. It enters the static Gaussian update for $\vect{\beta}$ through a
prior precision matrix. Let
$\mathcal{A}\subseteq\{1,\ldots,p\}$ denote the active shrinkage set; by
default an intercept, if present, is not shrunk. For $j\notin\mathcal{A}$ the
package uses a fixed intercept precision. The direct \code{rhs} path is not
developed here because its log-scale block is nonconjugate; the
\code{rhs_ns} path is the conjugate sparse-prior formulation documented in
this supplement.

\subsection{RHS-NS prior hierarchy}

For $j\in\mathcal{A}$, the \code{rhs_ns} hierarchy is
\begin{align}
\beta_j\mid\tau^2,\lambda_j^2,\zeta^2
&\propto
\N(\beta_j\mid0,\tau^2\lambda_j^2)
\N(\beta_j\mid0,\zeta^2),\\
\lambda_j^2\mid\nu_j &\sim \IG(1/2,1/\nu_j),&
\nu_j&\sim\IG(1/2,1),\\
\tau^2\mid\xi &\sim\IG(1/2,1/\xi),&
\xi&\sim\IG(1/2,1/\tau_0^2),\\
\zeta^2&\sim\IG(a_\zeta,b_\zeta),
\label{eq:rhs_ns_hierarchy}
\end{align}
unless $\zeta^2$ is fixed by the user. The resulting prior precision
contribution for active coefficients is
\begin{equation}
\pi_j=\frac{1}{\tau^2\lambda_j^2}+\frac{1}{\zeta^2}.
\label{eq:rhs_ns_precision}
\end{equation}
Consequently, the $\vect{\beta}$ update in MCMC or VB is obtained from the
static Gaussian update after replacing $\mat{B}_0^{-1}$ by a diagonal precision
matrix with active entries $\pi_j$ or their variational expectations.

\subsection{RHS-NS MCMC updates}

Conditioning on $\vect{\beta}$, the shrinkage block has closed-form
inverse-gamma updates:
\begin{align}
\lambda_j^2\mid\cdot
&\sim
\IG\left(1,\ \frac{1}{\nu_j}+\frac{\beta_j^2}{2\tau^2}\right),\\
\nu_j\mid\cdot
&\sim
\IG\left(1,\ 1+\frac{1}{\lambda_j^2}\right),\\
\tau^2\mid\cdot
&\sim
\IG\left(\frac{|\mathcal{A}|+1}{2},\
\frac{1}{\xi}+\frac{1}{2}\sum_{j\in\mathcal{A}}\frac{\beta_j^2}{\lambda_j^2}
\right),\\
\xi\mid\cdot
&\sim
\IG\left(1,\ \frac{1}{\tau_0^2}+\frac{1}{\tau^2}\right),\\
\zeta^2\mid\cdot
&\sim
\IG\left(a_\zeta+\frac{|\mathcal{A}|}{2},\
b_\zeta+\frac{1}{2}\sum_{j\in\mathcal{A}}\beta_j^2
\right),
\label{eq:rhs_ns_mcmc}
\end{align}
with the last line omitted when $\zeta^2$ is fixed.
These equations are the conditional targets for the RHS-NS scale block. The
implementation may temporarily freeze or schedule the global-scale update during
configured warmup iterations for numerical stability; this changes only the
update schedule, not the conditional target displayed above.

\subsection{RHS-NS VB updates and ELBO contribution}

The package uses a mean-field approximation over the RHS-NS scale block. For
$j\in\mathcal{A}$,
\begin{align}
q(\lambda_j^2)&=\IG(a_{\lambda j},b_{\lambda j}),&
q(\nu_j)&=\IG(a_{\nu j},b_{\nu j}),\\
q(\tau^2)&=\IG(a_\tau,b_\tau),&
q(\xi)&=\IG(a_\xi,b_\xi),\\
q(\zeta^2)&=\IG(a_\zeta^\star,b_\zeta^\star),
\end{align}
with $a_{\lambda j}=a_{\nu j}=a_\xi=1$ and
$a_\tau=(|\mathcal{A}|+1)/2$. Let
$\bar\beta_j^2=\E_q(\beta_j^2)$. The coordinate updates are
\begin{align}
b_{\lambda j}
&=\frac{1}{2}\bar\beta_j^2\,\E(\tau^{-2})+\E(\nu_j^{-1}),\\
b_{\nu j}
&=1+\E\{(\lambda_j^2)^{-1}\},\\
b_\tau
&=\frac{1}{2}\sum_{j\in\mathcal{A}}\bar\beta_j^2
\E\{(\lambda_j^2)^{-1}\}+\E(\xi^{-1}),\\
b_\xi
&=\frac{1}{\tau_0^2}+\E(\tau^{-2}),\\
a_\zeta^\star
&=a_\zeta+\frac{|\mathcal{A}|}{2},\qquad
b_\zeta^\star=b_\zeta+\frac{1}{2}\sum_{j\in\mathcal{A}}\bar\beta_j^2.
\label{eq:rhs_ns_vb_updates}
\end{align}
The third line is shorthand for
$b_\tau=\frac{1}{2}\sum_{j\in\mathcal{A}}\bar\beta_j^2
\E\{(\lambda_j^2)^{-1}\}+\E(\xi^{-1})$, where the expectation of
$(\lambda_j^2)^{-1}$ is component-specific inside the summation.
As in MCMC, optional warmup scheduling for the $\tau^2,\xi$ block affects when
the coordinate is updated, not the coordinate target itself.

The following term-by-term ELBO contribution matches the package's
\code{rhs_ns} scale block. For $x\sim\IG(a,b)$ under the rate
parameterization used above, define
\begin{align}
m_{\log x}(a,b)&=\log b-\psi(a),&
m_{x^{-1}}(a,b)&=a/b,\\
h_{\IG}(a,b)&=a+\log b+\log\Gamma(a)-(a+1)\psi(a),
\end{align}
where $\psi(\cdot)$ is the digamma function and $h_{\IG}$ is the entropy of an
inverse-gamma variational factor. Let
\begin{align}
\ell_{\lambda j}&=m_{\log x}(a_{\lambda j},b_{\lambda j}),&
r_{\lambda j}&=m_{x^{-1}}(a_{\lambda j},b_{\lambda j}),\\
\ell_{\nu j}&=m_{\log x}(a_{\nu j},b_{\nu j}),&
r_{\nu j}&=m_{x^{-1}}(a_{\nu j},b_{\nu j}),\\
\ell_\tau&=m_{\log x}(a_\tau,b_\tau),&
r_\tau&=m_{x^{-1}}(a_\tau,b_\tau),\\
\ell_\xi&=m_{\log x}(a_\xi,b_\xi),&
r_\xi&=m_{x^{-1}}(a_\xi,b_\xi).
\end{align}
If $\zeta^2$ is random, set
$\ell_\zeta=m_{\log x}(a_\zeta^\star,b_\zeta^\star)$ and
$r_\zeta=m_{x^{-1}}(a_\zeta^\star,b_\zeta^\star)$; if it is fixed at
$\zeta_0^2$, set $\ell_\zeta=\log\zeta_0^2$ and
$r_\zeta=1/\zeta_0^2$.

The coefficient-prior contribution is
\begin{align}
\elbo_{\beta,\mathrm{hs}}
&=
\sum_{j\in\mathcal{A}}
\left\{
-\frac12\log(2\pi)
-\frac12\ell_\tau
-\frac12\ell_{\lambda j}
-\frac12\bar\beta_j^2 r_\tau r_{\lambda j}
\right\},\\
\elbo_{\beta,\mathrm{slab}}
&=
\sum_{j\in\mathcal{A}}
\left\{
-\frac12\log(2\pi)
-\frac12\ell_\zeta
-\frac12\bar\beta_j^2 r_\zeta
\right\}.
\end{align}
The local, global, and slab prior contributions are
\begin{align}
\elbo_{\lambda\mid\nu}
&=
\sum_{j\in\mathcal{A}}
\left\{
-\frac12\ell_{\nu j}
-\log\Gamma(1/2)
-\frac32\ell_{\lambda j}
-r_{\nu j}r_{\lambda j}
\right\},\\
\elbo_{\nu}
&=
\sum_{j\in\mathcal{A}}
\left\{
-\log\Gamma(1/2)
-\frac32\ell_{\nu j}
-r_{\nu j}
\right\},\\
\elbo_{\tau\mid\xi}
&=
-\frac12\ell_\xi
-\log\Gamma(1/2)
-\frac32\ell_\tau
-r_\xi r_\tau,\\
\elbo_{\xi}
&=
\frac12\log(\tau_0^{-2})
-\log\Gamma(1/2)
-\frac32\ell_\xi
-\tau_0^{-2}r_\xi,\\
\elbo_{\zeta}
&=
\begin{cases}
a_\zeta\log b_\zeta-\log\Gamma(a_\zeta)
-(a_\zeta+1)\ell_\zeta-b_\zeta r_\zeta,
& \zeta^2\text{ random},\\
0, & \zeta^2\text{ fixed}.
\end{cases}
\end{align}
The entropy contribution is
\begin{equation}
\elbo_{\mathrm{ent}}
=
\sum_{j\in\mathcal{A}} h_{\IG}(a_{\lambda j},b_{\lambda j})
+\sum_{j\in\mathcal{A}} h_{\IG}(a_{\nu j},b_{\nu j})
+h_{\IG}(a_\tau,b_\tau)
+h_{\IG}(a_\xi,b_\xi)
+\Ind(\zeta^2\text{ random})h_{\IG}(a_\zeta^\star,b_\zeta^\star).
\end{equation}
Finally, when the intercept is not included in $\mathcal{A}$, the package adds
the fixed Gaussian intercept-prior contribution
\begin{equation}
\elbo_{\beta_0}
=
\frac12\{\log(\pi_0)-\log(2\pi)\}
-\frac12\pi_0\bar\beta_0^2,
\end{equation}
where $\pi_0$ is the fixed intercept precision. If the intercept is shrunk,
this term is omitted because the intercept is already represented in the
active RHS-NS block.

Thus the package contribution is
\begin{equation}
\elbo_{\mathrm{rhs\_ns}}
=
\elbo_{\beta,\mathrm{hs}}
+\elbo_{\beta,\mathrm{slab}}
+\elbo_{\lambda\mid\nu}
+\elbo_{\nu}
+\elbo_{\tau\mid\xi}
+\elbo_{\xi}
+\elbo_{\zeta}
+\elbo_{\mathrm{ent}}
+\elbo_{\beta_0}.
\label{eq:rhs_ns_elbo}
\end{equation}
All terms are available in closed form because each scale factor is
inverse-gamma. This is the RHS-NS component of the full static-model ELBO; the
likelihood, latent-variable, and $(\sigma,\gamma)$ terms are added by the
surrounding static AL or exAL variational routine.

\section{Bounded slice sampling and Laplace-Delta blocks}
\label{sec:slice}

\subsection{Bounded univariate slice sampler}

Let $h(x)$ be an unnormalized log-density on a finite interval $[a,b]$, and let
$x_0\in(a,b)$. The bounded slice update implemented in the package is:
\begin{enumerate}[label=(\arabic*)]
\item Draw $e\sim\Exp(1)$ and set the log-slice height $z=h(x_0)-e$.
\item Draw $u\sim\operatorname{Unif}(0,w)$ and initialize
$L=x_0-u$, $R=x_0+w-u$.
\item Step out left and right by increments of $w$, respecting the bounds
$a$ and $b$, until the interval endpoints lie below the slice or the maximum
number of expansions is reached.
\item Truncate the interval to $[a,b]$.
\item Repeatedly draw $x_1\sim\operatorname{Unif}(L,R)$. If $h(x_1)\ge z$,
accept $x_1$. Otherwise shrink the interval to $[L,x_1]$ if $x_1>x_0$ and to
$[x_1,R]$ otherwise.
\end{enumerate}
For dynamic and static exAL MCMC fits, this sampler is applied to the
one-dimensional $\gamma$ target on $(L,U)$ by default. The same bounded-slice
utility can be reused by other internal one-dimensional nonconjugate blocks;
the user-facing algorithms covered here are the exAL updates described above.

\subsection{Laplace-Delta approximation}

For a nonconjugate block $\vect{\eta}\in\R^d$ with log target
$\ell(\vect{\eta})$, let $\hat{\vect{\eta}}$ be a local mode and
$\mat{H}=\nabla^2\ell(\hat{\vect{\eta}})$. If $-\mat{H}$ is positive definite,
the Laplace approximation is
\begin{equation}
q(\vect{\eta})\approx
\N\left(\hat{\vect{\eta}},\ [-\mat{H}]^{-1}\right).
\label{eq:ld_gaussian}
\end{equation}
For a smooth function $r(\vect{\eta})$, the Delta approximation used by LDVB is
\begin{equation}
\E\{r(\vect{\eta})\}
\approx
r(\hat{\vect{\eta}})
+
\frac{1}{2}\tr\left[
\nabla^2 r(\hat{\vect{\eta}})[-\mat{H}]^{-1}
\right].
\label{eq:delta_generic}
\end{equation}
This is used for smooth functions of $(\sigma,\gamma)$ in dynamic and static
exAL models.

\section{Package-to-algorithm map}

\begin{table}[!htb]
\centering
\begin{tabularx}{\textwidth}{@{}p{0.24\textwidth}>{\raggedright\arraybackslash}X@{}}
\toprule
Function & Main algorithmic role \\
\midrule
\code{exdqlmMCMC()} & Dynamic DQLM/exDQLM posterior simulation: FFBS for states, Gibbs for conjugate latent/scale blocks, and bounded slice for $\gamma$ when exAL is used. \\
\code{exdqlmLDVB()} & Dynamic DQLM/exDQLM mean-field VB: CAVI for Gaussian/GIG/truncated-normal blocks and Laplace-Delta for $(\sigma,\gamma)$ in exAL. \\
\code{exalStaticMCMC()} & Static AL/exAL posterior simulation: Gaussian coefficient update, Gibbs latent/scale updates, RHS-NS inverse-gamma scale updates when selected, and bounded slice for $\gamma$ in exAL. \\
\code{exalStaticLDVB()} & Static AL/exAL mean-field VB: Gaussian coefficient factor, CAVI latent and RHS-NS scale factors, and Laplace-Delta for $(\sigma,\gamma)$ in exAL. \\
\bottomrule
\end{tabularx}
\caption{Connection between the main \pkg{exdqlm} functions and the
algorithmic blocks described in this supplement.}
\label{tab:supp_algorithm_map}
\end{table}

\end{document}
